    \subsection{Calor Latente de Fusión.}
        Se calibró la báscula con la masa del calorímetro. Luego se vertió agua en el calorímetro, se regristró su masa con ayuda de la báscula electrónica , así como su temperatura. Se vertió el hielo en el calorímetro y se dejó disolver en el agua, posteriormente se registró la temperatura de equilirio entre el hielo y el agua.
    
    \subsection{Coeficiente de Dilatación.}
        Se midió la longitud de los lingotes con un vernier, así como la altura y ancho con un tornillo micrométrico. con el tornillo se midió el diámetro del balín. Se guardó el balín y los lingotes en el refrigerador del  aboratorio durante 20 minutos. Luego se registró de nuevo el largo, ancho y diámetro del balín. 

